% ═══════════════════════════════════════════════════════════════════
\subsection{Résultats quantitatifs}
\label{subsec:ch3-resultats-quantitatifs}
% ═══════════════════════════════════════════════════════════════════

Le protocole décrit à la section précédente a été appliqué sur la
partition de test de CBIS-DDSM, composée de \textbf{787~cas} issus de
\textbf{349~patientes} distinctes (475~bénins, 312~malins). Aucune de
ces patientes n'apparaît dans les ensembles d'entraînement ou de
validation, garantissant ainsi l'absence de fuite d'information.

\paragraph{Métriques globales.}
Le tableau~\ref{tab:ch3-metriques-test} rassemble les indicateurs
obtenus avec le modèle final (\textit{SWA}, étape~2) et un
\textit{test-time augmentation} à 16~vues (TTA-16). La
\textbf{spécificité} atteint 76,2\,\%, ce qui traduit une capacité
satisfaisante à écarter les lésions bénignes et à limiter le nombre de
biopsies inutiles. La \textbf{sensibilité} de 63,1\,\% reste en deçà
du seuil clinique souhaitable~: dans un contexte de dépistage, il est
préférable de maximiser la détection des cas malins, quitte à accepter
davantage de faux positifs. Cette limitation est discutée plus bas.

\begin{table}[H]
    \centering
    \caption{Performances du classificateur EfficientNet-B3 sur le test
             CBIS-DDSM (TTA-16, seuil $\tau = 0{,}535$).}
    \label{tab:ch3-metriques-test}
    \begin{tabular}{ll}
        \toprule
        \textbf{Métrique} & \textbf{Valeur} \\
        \midrule
        AUC-ROC            & 0,7812 \\
        AUC-PR             & 0,7183 \\
        Sensibilité        & 63,1\,\% \\
        Spécificité        & 76,2\,\% \\
        Exactitude          & 71,0\,\% \\
        Score F1            & 0,6334 \\
        MCC                 & 0,394 \\
        \bottomrule
    \end{tabular}
\end{table}

\paragraph{Matrice de confusion.}
La figure~\ref{fig:ch3-roc-cm} présente, à gauche, la courbe
\textit{ROC} du modèle avec le point de fonctionnement retenu
($\tau = 0{,}535$), et à droite la matrice de confusion correspondante.
Sur les 312~cas malins, le classifieur en identifie correctement
197~(\textbf{vrais positifs}), tandis que 115~échappent à la détection
(\textbf{faux négatifs}). Côté bénin, 362~des 475~cas sont correctement
classés (\textbf{vrais négatifs}), et 113~sont orientés à tort vers une
investigation complémentaire (\textbf{faux positifs}).

\begin{figure}[H]
    \centering
    \includegraphics[width=0.95\textwidth]{images/cls_roc_cm_final.png}
    \caption[Courbe ROC et matrice de confusion du classificateur
             EfficientNet-B3]{Courbe \textit{ROC}
             (AUC~=~0,7812) et matrice de confusion
             ($\tau = 0{,}535$) obtenus sur le test CBIS-DDSM
             (787~cas, TTA-16). Le point jaune indique le seuil
             de décision retenu.}
    \label{fig:ch3-roc-cm}
\end{figure}

\paragraph{Discussion clinique.}
En mammographie de dépistage, un \textbf{faux négatif} constitue
l'erreur la plus coûteuse~: une lésion maligne non détectée retarde le
diagnostic et peut compromettre le pronostic de la patiente. La
sensibilité observée (63,1\,\%) indique que le modèle, utilisé seul,
laisserait passer environ un tiers des cas malins. Ce résultat souligne
que le système proposé ne peut se substituer au radiologue, mais
s'inscrit comme un \textbf{outil d'aide à la décision}~: il signale les
régions suspectes identifiées par la segmentation et fournit une
probabilité de malignité que le praticien intègre à son analyse globale.

La valeur de l'\textbf{AUC} (0,7812) se situe dans la fourchette
rapportée par la littérature pour des classifieurs entraînés sur
CBIS-DDSM, généralement comprise entre 0,75 et
0,85~\cite{cbisddsm2017, shia2025mammography}. Il convient de noter que
CBIS-DDSM est un jeu de données particulièrement exigeant~: les
\textit{patches} sont recadrés à partir de mammographies complètes avec
des annotations parfois ambiguës, et le déséquilibre modéré entre
classes (59\,\% de bénins) complique la tâche de discrimination.

Le \textbf{score F1} de 0,6334 et le coefficient de corrélation de
Matthews (\textbf{MCC}~=~0,394) confirment une performance supérieure au
hasard, tout en révélant une marge d'amélioration. Parmi les pistes
envisageables figurent l'augmentation de la taille du jeu
d'entraînement, l'intégration d'informations contextuelles (vue
controlatérale, antécédents) ou le recours à des architectures plus
récentes. Ces perspectives sont développées en conclusion du chapitre.
